% Part: proof-theory
% Chapter: normalization
% Section: normalization-thm

\documentclass[../../../include/open-logic-section]{subfiles}

\begin{document}

\olfileid[pt]{pt}{nor}{thm}

\olsection{O Teorema da Normalização}

!!^a{proof} está na forma \emph{normal} se não tem segmentos de corte.
Claramente, esse é o caso se, e somente se, não podemos aplicar a ela nem uma
conversão de permutação nem uma conversão de redução. \emph{Normalizar}
!!a{proof} significa transformá-la em uma prova normal aplicando
conversões de permutação e de redução até que não reste nenhum corte. O fato de
isso ser sempre possível é o \emph{teorema da normalização}.

\begin{thm}\ollabel{thm:normalization} Toda \Log{N1i}-!!{proof}
$\delta$ se normaliza, isto é, pode ser transformada em
!!a{proof}~$\delta'$ sem cortes.
\end{thm}

\begin{proof}
  Por indução sobre o posto e o comprimento de corte de~$\delta$. A hipótese de
  indução é que qualquer !!{proof} que tenha um posto de corte menor, ou
  o mesmo posto de corte mas um comprimento de corte menor, se normaliza.

  Se o comprimento de corte for~$0$, não há cortes e $\delta$ já é
  normal. Portanto, suponha que o posto e o comprimento de corte de~$\delta$ sejam~$>
  0$. Escolha o segmento de corte maximal mais acima e mais à direita. Se o
  segmento de corte tiver comprimento $>1$, aplique uma conversão de permutação. Se o
  segmento de corte tiver comprimento~$1$, aplique a conversão de redução
  correspondente. Isso resulta em uma nova !!{proof} que tem ou o
  mesmo posto de corte e um comprimento de corte menor, ou um posto de corte menor; então
  a hipótese de indução se aplica.
\end{proof}

A estratégia usada na demonstração acima (sempre tratar de um
corte maximal mais acima e mais à direita) sugere uma ordem específica em que
as conversões devem ser aplicadas a !!a{proof} para resultar em uma
!!{proof} normal. Surpreendentemente, a ordem na verdade não
importa. Aplicar conversões de permutação e de redução em qualquer ordem
acaba resultando em uma prova normal.

A definição de segmento e de corte pode ser formulada facilmente para
\Log{N2i}, e as conversões de permutação e de redução para \Log{N1i}
se traduzem diretamente para \Log{N2i}. Consequentemente, o teorema da
normalização também vale para \Log{N2i}.

\begin{thm}\ollabel{thm:normalization-N2i} Toda \Log{N2i}-!!{proof}
$\delta$ se normaliza, isto é, pode ser transformada em
!!a{proof}~$\delta'$ sem cortes.
\end{thm}

\end{document}
